\chapter{Introdução}

Parte inicial do texto, onde devem constar a delimitação do assunto tratado, objetivos da pesquisa e outros elementos necessários para situar o tema do trabalho. Uma dica: quando a sigla aparece pela primeira vez no texto, a forma completa do seu significado a precede, sendo a sigla colocada entre parênteses. Por exemplo: Associação Brasileira de Normas Técnicas (ABNT). A introdução  deve preencher de duas a três páginas inteiras, e deve enfatizar a motivação para o trabalho, o estado da arte. os objetivos e, ao final, um resumo de como a dissertação está organizada, quantos capítulos e o que será tratado em cada capítulo.

Não esquecer de referenciar todo o texto escrito no documento. Para isso, pode ser utilizado o sistema numérico de referências (ex: [1]), ou o sistema autor-data \cite{NBR6023}, conforme mostrado mais adiante.

Este modelo em \LaTeX{} foi feito de forma a reproduzir ao máximo a formatação do Modelo de Dissertação em docx, disponibilizado no Portal do PPGEE\footnote{Disponível em: \url{https://www.ifpb.edu.br/ppgee/documentos/modelo-de-dissertacao/modelo_dissertacao_ppgee.docx}. Acesso em 30 de dezembro de 2021.}. A classe \texttt{ppgee.cls} consiste em adaptações e configurações da classe \texttt{abntex2}. Consulte o manual do pacote \texttt{abntex2} \cite{abntex2classe} e o seu modelo de trabalho acadêmico \cite{abntex2modelo} para mais informações e possibilidades que não estejam exemplificadas neste documento.

O \LaTeX, juntamente com a classe \texttt{ppgee.cls} e o pacote \texttt{abntex2} já cuidam da formatação de fontes, parágrafos, figuras e tabelas, de forma que é preciso apenas se ocupar com o conteúdo. A numeração de figuras, tabelas e equações também não necessita de nenhuma intervenção, sendo realizada de forma automática pelo \LaTeX.

É possível escolher entre as fontes Arial ou Times New Roman. Por padrão, a classe utiliza a fonte Arial. Se quiser utilizar a fonte Times New Roman, adicione a opção \texttt{fonte=times} nas opções da classe.